% Figure: Cantor diagonalisation --- flip each marked digit, escape the list.
% Source for Figure~\ref{fig:diagonal}. TikZ only; regenerable by compiling main.tex.
\begin{tikzpicture}[
  dg/.style={font=\footnotesize\ttfamily,text=ink},
  hit/.style={draw=amber,thick,rounded corners=1pt},
  lab/.style={font=\footnotesize\sffamily,text=slate},
  new/.style={font=\footnotesize\ttfamily,text=accent},
]
\def\cw{0.21}
% the table (digits monospaced); rows 0..7, decimal digits at columns 2..11
\foreach \r/\row in {
  0/{0.5000000000},
  1/{0.3333333333},
  2/{0.1415926535},
  3/{0.7182818284},
  4/{0.4142135623},
  5/{0.6180339887},
  6/{0.1010010001},
  7/{0.2500000000}}
  \node[dg,anchor=west] at (0,-\r*0.34) {\row};
% amber boxes on the diagonal digit of each row: column index = 2 + r
\foreach \r in {0,...,7}
  \node[hit,minimum width=2.6mm,minimum height=3mm] at ({(2+\r)*\cw + 0.018},-\r*0.34-0.02) {};
\node[lab,anchor=west] at (2.7,-0.1) {each row's marked digit};
\draw[->,slate] (2.75,-0.45) -- (1.28,-1.35);
% constructed row
\node[new,anchor=west] at (0,-2.95) {0.45555555...};
\node[lab,anchor=west] at (2.2,-2.98) {flipped each marked digit --- absent from every row};
\end{tikzpicture}
